\section{Canonical deeper transient chains}
\label{sec:deep}

We now iterate the one-level mechanism in one specific way.  Fix $L\geq1$
and add transient states $r_0,\ldots,r_{L-1}$.  The only transient edges are
\[
 r_0\longrightarrow r_1\longrightarrow\cdots\longrightarrow r_{L-1}
 \longrightarrow\bigcup_{s=0}^{p-1}V_s.
\]
An arrow to a single transient state forces all $d$ children to carry that
state.  The final state $r_{L-1}$ sees every core label, and the core has no
edge back to the chain.  We call this the \emph{canonical unrestricted
$L$-level forced chain}.  Starting at $r_0$, level $L$ consists of
\[
 N=d^L
\]
ordered core roots whose phases may be selected independently.

For $m\in\Comp_p(d^L)$, define
\begin{equation}
 D_L(m;c)=\min_j\frac1{d^L}\sum_{s=0}^{p-1}m_sH_{s+j}(c),
 \qquad
 D_L^*(c)=\max_{m\in\Comp_p(d^L)}D_L(m;c).
 \label{eq:DL-definition}
\end{equation}

\begin{theorem}[Exact forced-chain optimizer]
\label{thm:deep-exact}
The full Markov hom tree-shift consisting of the complete cyclic core and the
canonical unrestricted $L$-level forced chain has Hausdorff dimension
\begin{equation}
 \dimH T_{M_L(a)}=D_L^*(c).
 \label{eq:deep-full-dimension}
\end{equation}
For a fixed phase assignment at level $L$ with composition $m$, its top-root
stratum has dimension $D_L(m;c)$.  In particular, the denominator in
\eqref{eq:DL-definition} is exactly $d^L$.
\end{theorem}

\begin{proof}
Fix a phase assignment of the $d^L$ level-$L$ core roots.  At total depth
$n\geq L$, its exact cylinder count is
\begin{equation}
 P_{n,L,m}=\prod_{s=0}^{p-1}\prod_{\ell=0}^{n-L}
 a_{s+\ell}^{m_s d^\ell}.
 \label{eq:deep-prefix-count}
\end{equation}
The $L$ transient levels have forced labels and contribute no multiplicity.
Since
\begin{equation}
 \frac{|\Delta_{n-L}|}{|\Delta_n|}\longrightarrow d^{-L},
 \label{eq:deep-scale-ratio}
\end{equation}
the proof of \cref{prop:fixed-one-level}, now using
\cref{lem:weighted-limit}, gives the residue limits
\[
 \frac1{d^L}\sum_sm_sH_{s+j}(c).
\]
The cylinder lemma proves the fixed-composition statement.  There are only
$p^{d^L}$ phase assignments, so their finite union has dimension $D_L^*(c)$
at the top root.

The full shift also allows roots in the core or at $r_q$ for $q>0$.  A root
at $r_q$ has $K=L-q$ transient levels and therefore value $D_K^*(c)$.  If
$m\in\Comp_p(d^K)$, then $dm\in\Comp_p(d^{K+1})$ and
\begin{equation}
 D_{K+1}(dm;c)=D_K(m;c).
 \label{eq:grid-embedding}
\end{equation}
Thus $D_{K+1}^*(c)\geq D_K^*(c)$.  A composition concentrated in one phase
also recovers $\min_jH_j(c)$ at every $K$, so the top-root value dominates
the core-root strata.  The finite-union identity now proves
\eqref{eq:deep-full-dimension}.
\end{proof}

\begin{theorem}[Finite-depth saturation and convergence]
\label{thm:deep-convergence}
For the canonical unrestricted chain,
\begin{align}
 D_L(m;c)=\bar c
 &\quad\Longleftrightarrow\quad
 \sum_sm_sc_{s+k}\text{ is independent of }k,
 \label{eq:DL-saturation}\\
 D_L^*(c)&\leq\bar c,
 \label{eq:DL-upper}\\
 D_{L+1}^*(c)&\geq D_L^*(c),
 \label{eq:DL-monotone}\\
 0\leq\bar c-D_L^*(c)
 &\leq\frac{p\max_jH_j(c)}{d^L}.
 \label{eq:DL-rate}
\end{align}
Consequently $D_L^*(c)\to\bar c$.  If $p\mid d^L$, equality already holds
at level $L$.  Under the full nonzero Fourier-support hypothesis of
\cref{thm:fourier-divisibility}, a saturating level-$L$ composition exists
only if $p\mid d^L$; no such converse is asserted without that hypothesis.
\end{theorem}

\begin{proof}
The quantity $D_L(m;c)$ is $\Psi_*^{(N)}(m;c)$ from
\eqref{eq:Psi-definition} with $N=d^L$.  Therefore
\eqref{eq:DL-saturation} and \eqref{eq:DL-upper} follow from
\cref{thm:saturation}, and the divisibility statements follow from
\cref{thm:fourier-divisibility}.  Monotonicity is exactly the embedding
\eqref{eq:grid-embedding}.

For the rate, choose a balanced composition $m\in\Comp_p(N)$, $N=d^L$,
whose entries are $\lfloor N/p\rfloor$ or $\lceil N/p\rceil$.  Write
$m_s=N/p+e_s$.  Then $\sum_se_s=0$ and $|e_s|<1$.  Mean preservation gives,
for every $j$,
\begin{align*}
 \left|\frac1N\sum_sm_sH_{s+j}(c)-\bar c\right|
 &=\left|\frac1N\sum_se_sH_{s+j}(c)\right|\\
 &\leq\frac{p\max_qH_q(c)}{N}.
\end{align*}
The optimized minimum is at least the minimum produced by this balanced
composition and at most its phase mean $\bar c$.  This proves
\eqref{eq:DL-rate} and the limit.
\end{proof}

\begin{figure}[t]
  \centering
  \begin{tikzpicture}
\begin{groupplot}[
  group style={group size=2 by 1,horizontal sep=1.35cm},
  width=0.47\textwidth,
  height=0.36\textwidth,
  grid=major,
  grid style={black!8},
  tick label style={font=\scriptsize},
  label style={font=\small},
  legend style={font=\scriptsize,draw=none,fill=none},
]
\nextgroupplot[
  xlabel={transient depth $L$},
  ylabel={dimension},
  xtick={1,2,3},
  legend style={at={(0.02,0.98)},anchor=north west,font=\tiny,draw=none,fill=white,fill opacity=0.85,text opacity=1},
]
\addplot+[blue,thick,mark=*] table[x=level,y=dimension,col sep=comma,
  restrict expr to domain={\thisrow{series}}{0:0}] {../figures/data/fig2_exact.csv};
\addlegendentry{$d=2,p=3,a=(1,2,3)$}
\addplot+[orange!85!black,thick,mark=square*] table[x=level,y=dimension,col sep=comma,
  restrict expr to domain={\thisrow{series}}{1:1}] {../figures/data/fig2_exact.csv};
\addlegendentry{$d=2,p=4,a=(1,2,3,1)$}
\addplot+[green!50!black,thick,mark=triangle*] table[x=level,y=dimension,col sep=comma,
  restrict expr to domain={\thisrow{series}}{2:2}] {../figures/data/fig2_exact.csv};
\addlegendentry{$d=3,p=2,a=(1,2)$}
\addplot+[blue,densely dashed,forget plot] coordinates {(1,0.5972531564093517) (3,0.5972531564093517)};
\addplot+[orange!85!black,densely dashed,forget plot] coordinates {(1,0.4479398673070138) (3,0.4479398673070138)};
\addplot+[green!50!black,densely dashed,forget plot] coordinates {(1,0.3465735902799727) (3,0.3465735902799727)};

\nextgroupplot[
  xlabel={transient depth $L$},
  ylabel={gap to $\bar c$},
  ymode=log,
  xtick={1,2,3,4,5,6,7,8},
  legend style={at={(0.02,0.02)},anchor=south west,font=\tiny,draw=none,fill=white,fill opacity=0.85,text opacity=1},
]
\addplot+[blue,thick,mark=*] table[x=level,y=balanced_gap,col sep=comma,
  restrict expr to domain={\thisrow{series}}{0:0}] {../figures/data/fig2_balanced.csv};
\addlegendentry{gap $(d,p)=(2,3)$}
\addplot+[blue,densely dashed,mark=none] table[x=level,y=certificate,col sep=comma,
  restrict expr to domain={\thisrow{series}}{0:0}] {../figures/data/fig2_balanced.csv};
\addlegendentry{bound $(2,3)$}
\addplot+[orange!85!black,thick,mark=square*] table[x=level,y=balanced_gap,col sep=comma,
  restrict expr to domain={\thisrow{series}}{1:1}] {../figures/data/fig2_balanced.csv};
\addlegendentry{gap $(d,p)=(3,4)$}
\addplot+[orange!85!black,densely dashed,mark=none] table[x=level,y=certificate,col sep=comma,
  restrict expr to domain={\thisrow{series}}{1:1}] {../figures/data/fig2_balanced.csv};
\addlegendentry{bound $(3,4)$}
\end{groupplot}
\end{tikzpicture}

  \caption{Finite-depth phase allocation.  \emph{Left:} exhaustive exact
  optimizers from the frozen composition sweep for three profiles at
  $L=1,2,3$; dashed horizontal lines are the corresponding spectral means.
  The $d=2,p=4$ profile saturates when $4\mid2^L$.  \emph{Right:} gaps of the
  displayed balanced compositions for $a=(2,3,4)$ with $(d,p)=(2,3)$ and
  $a=(2,3,4,2)$ with $(d,p)=(3,4)$, together with the proved certificate
  $p\max_jH_j(c)/d^L$.  The right panel shows a certificate, not an assertion
  that the balanced composition is always the optimizer.}
  \label{fig:deep-convergence}
\end{figure}

\begin{proposition}[A bounded restricted extension]
\label{prop:balanced-access}
Let $F_L\subseteq\Comp_p(d^L)$ be a nonempty finite family of compositions
admitted by a separately declared depth-$L$ feeder, with the same complete
core and no return edge.  Its top-root dimension is
$\max_{m\in F_L}D_L(m;c)$.  If there is a constant $C$ and a choice
$m^{(L)}\in F_L$ for every $L$ such that
\begin{equation}
 \max_s\left|m_s^{(L)}-\frac{d^L}{p}\right|\leq C,
 \label{eq:balanced-access}
\end{equation}
then
\begin{equation}
 0\leq \bar c-\max_{m\in F_L}D_L(m;c)
 \leq\frac{Cp\max_jH_j(c)}{d^L}.
 \label{eq:restricted-rate}
\end{equation}
Monotonicity for restricted families additionally requires a nesting
condition such as $dF_L\subseteq F_{L+1}$; it is not inferred from
\eqref{eq:balanced-access} alone.
\end{proposition}

\begin{proof}
The finite-union argument in \cref{thm:deep-exact} gives the exact maximum
over $F_L$.  Write $m_s^{(L)}=d^L/p+e_s$ and repeat the last estimate in the
proof of \cref{thm:deep-convergence}, now with $|e_s|\leq C$.
\end{proof}

The canonical chain is therefore more than a continuous-simplex heuristic:
each finite $L$ is an exact integer optimization on the grid with denominator
$d^L$.  The limit follows because these grids contain balanced points with
uniformly bounded coordinate discrepancy.
